%% kapa: exam=1 round=20 year=2009 period=2 form=A subject=accounting subject_ko=회계학 seq=34 number=34
%% source: Q-Net 감정평가사 기출문제 (cst00309) artlSeq=5024178 file=extracted/1008/5024178/2009년 제20회 감정평가사 1차시험 기출문제/2교시.hwp
%% conversion: kordoc 4.24.51 -> tools/gen1.py (auto); review=proofread-sonnet; flags=none
\begin{question}{34}
(주)한국은 A타입과 B타입 두 종류의 제품을 생산한다. 제품별 생산량 및 단위당 재료원가와 노무원가는 다음과 같다.

{\small\centering
\begin{tabular}{@{}l r r@{}}
\hline
 & \multicolumn{1}{c}{\uline{A타입}} & \multicolumn{1}{c}{\uline{B타입}} \\
생산량 & 4,000단위 & 5,000단위 \\
단위당 직접재료원가 & ￦200 & ￦250 \\
단위당 직접노무원가 & ￦10 & ￦20 \\
\hline
\end{tabular}\par}
회사의 제조활동은 다음 3가지로 구분되며, 활동별 제조간접원가와 관련된 자료는 다음과 같다.

{\small\centering
\begin{tabular}{@{}c c r r r r@{}}
\hline
 & & & \multicolumn{3}{c}{제품별 원가동인 수} \\
\cline{4-6}
\multirow{2}{*}{활동} & \multirow{2}{*}{\makecell{원가\\동인}} & \multirow{2}{*}{\makecell{활동별\\제조간접원가}} & \makecell{A타입\\(4,000단위)} & \makecell{B타입\\(5,000단위)} & \multicolumn{1}{c}{계} \\
 & & & & & \\
\cline{1-6}
\makecell{구매\\주문} & \makecell{주문\\횟수} & ￦160,000 & 15회 & 25회 & 40회 \\
\makecell{품질\\검사} & \makecell{검사\\횟수} & 570,000 & 10회 & 20회 & 30회 \\
\makecell{기계\\가동} & \makecell{기계\\시간} & \uline{4,000,000} & 80시간 & 120시간 & 200시간 \\
 & & ￦\uline{4,730,000} & & & \\
\hline
\end{tabular}\par}
활동기준원가계산(ABC)에 의하여 A타입과 B타입의 단위당 제조원가를 계산하면 각각 얼마인가? (단, 계산시 소수점이하 첫째자리에서 반올림한다.) \par

{\small\centering
\begin{tabular}{@{}ccc@{}}
 & A타입 & B타입 \\
\cline{2-2}\cline{3-3}
① & ￦485 & ￦325 \\
② & ￦485 & ￦898 \\
③ & ￦495 & ￦846 \\
④ & ￦673 & ￦846 \\
⑤ & ￦673 & ￦898 \\
\end{tabular}\par}
\end{question}
\answer{4}
